Для замыкания дискретной задачи необходимо задать граничные условия на давление~$p$ и (при использовании модели JFO) на степень заполнения~$\vartheta$. Тип граничных условий определяется физической постановкой (подразделы~2.1.2-2.1.4) и геометрией расчётной области. Ниже систематизированы три основных типа граничных условий для давления.

\textit{Дирихле (заданное давление).}
На соответствующем участке границы $\Gamma_D$ давление фиксируется:
\[
  p = p_0 \quad \text{на} \quad \Gamma_D.
\]

Наиболее распространённый случай -- $p_0 = p_{\mathrm{cav}} = 0$ (в избыточных давлениях) на открытых границах расчётной области, где смазочный слой сообщается с окружающей средой. Типичные примеры: торцы по оси~$z$ для радиального подшипника, радиальные границы $r = R_{\mathrm{in}},\, R_{\mathrm{out}}$ для упорного.

\textit{Периодичность.}
По окружному (периодическому) направлению давление совпадает на противоположных границах области:
\[
  p(x_{1,s},\, x_2) = p(x_{1,t},\, x_2).
\]

Данное условие применяется в упорных и радиальных подшипниках, где расчётная область охватывает полный период по окружной координате. При расчёте сектора вместо периодичности на границах сектора задаются условия симметрии ($\partial p / \partial n = 0$) или заданное давление в зависимости от постановки.

\textit{Непроницаемость (Нейман).}
На участках с условием симметрии или непроницаемости нормальная производная давления обращается в нуль:
\[
  \frac{\partial p}{\partial n} = 0 \quad \text{на} \quad \Gamma_N.
\]

Условие $\partial p / \partial n = 0$ эквивалентно отсутствию напорного потока через границу и используется на линиях симметрии или на искусственных границах при расчёте сектора.

В таблице~\ref{tab:bc_summary} сведены граничные условия для давления, соответствующие трём постановкам, сформулированным в подразделах~2.1.2-2.1.4.

\begin{table}[!ht]
\centering
\caption{Граничные условия для давления в различных постановках
  (координаты $x_1$,~$x_2$ определены в подразделах~2.1.6 и~2.2.1)}
\label{tab:bc_summary}
\begin{tabularx}{\textwidth}{|X|c|c|}
\hline
\multicolumn{1}{|c|}{\textbf{Постановка}} & \textbf{По $\boldsymbol{x_1}$}
                       & \textbf{По $\boldsymbol{x_2}$} \\
\hline
Торцевой зазор
  & периодичность
  & $p = p_{\mathrm{cav}}$, $x_2 = r$ \\
\hline
Радиальный зазор при окружном сдвиге
  & периодичность
  & $p = p_{\mathrm{cav}}$, $x_2 = z$ \\
\hline
Радиальный зазор при возвратно-поступательном движении
  & $p = p_{\mathrm{cav}}$, $x_1 = z$
  & периодичность \\
\hline
\end{tabularx}
\end{table}

\FloatBarrier

При использовании массо-сохраняющей модели кавитации JFO (раздел~2.3) дополнительно задаются граничные условия на степень заполнения~$\vartheta$. На границах подачи смазки полагается $\vartheta = 1$ (полное заполнение зазора). По периодическим направлениям степень заполнения периодична, аналогично давлению. На выходных границах с $p = p_{\mathrm{cav}}$ величина~$\vartheta$ явно не фиксируется; она определяется из уравнения сохранения~\eqref{eq:jfo_conservation} и схемы переноса, что соответствует свободному выходу смеси.

Во всех рассматриваемых постановках кавитационное давление принимается равным нулю в избыточных давлениях: $p_{\mathrm{cav}} = 0$.

Перечисленные типы граничных условий по-разному реализуются
в сеточной схеме. Узлы с условием Дирихле исключаются
из итерационного обновления: их значения фиксируются
и при вычислении потоков на соседних гранях используются
как известные. Периодичность по координатe~$s_1$ реализуется
отождествлением узлов на противоположных границах расчётной
области~-- фактически граница «замыкается» в кольцо,
и узел $i = 0$ совпадает с узлом $i = N_1$. Условие
неотрицательности давления ($P \geq 0$, кавитация
Рейнольдса) обеспечивается операцией усечения после каждого
шага SOR: если обновлённое значение оказывается отрицательным,
оно принудительно полагается равным нулю. Это эквивалентно
упрощённому методу активного множества без пересчёта степени
заполнения~$\vartheta$. Для JFO-постановки трактовка
кавитационной границы принципиально иная: граница является
свободной и определяется из условия непрерывности потока;
её поточечное обновление описано в разделе~3.5.

\textit{Физический смысл периодичности и~её~отличие от~искусственных граничных условий.}
Условие периодичности по~окружной координате является не~искусственным техническим приёмом, а~точным отражением физической геометрии задачи. Для~радиального подшипника окружная координата~$\varphi$ пробегает полный период $[0, 2\pi]$, и~давление на~входе в~расчётную область действительно равно давлению на~выходе~-- обе~<<границы>> на~самом деле описывают одно и~то~же физическое сечение. Это отличает условие периодичности от~условий Дирихле или~Неймана, которые описывают реальные физические границы (открытые торцы, стенки). Нарушение условия периодичности (например, замена его~на~условие постоянного давления на~обеих граничных линиях~$\varphi = 0$ и~$\varphi = 2\pi$) привело бы к~существенно неверным результатам, так как~фактически изменило бы топологию задачи.

\textit{Граница кавитации в~JFO-постановке как~свободная граница.}
В~упрощённых моделях кавитации (обнуление давления, модель Гюмбеля) граница кавитации явно не~выделяется: давление просто обнуляется поточечно в~тех~узлах, где оно~оказывается отрицательным. Никакого специального граничного условия на~линии раздела зон не~задаётся~-- эта линия неявно определяется самим алгоритмом обнуления.

В~JFO-постановке граница кавитации является свободной границей в~математическом смысле: её~положение не~фиксируется заранее, а~определяется в~ходе решения как~часть ответа. На~этой границе должны выполняться одновременно два условия: давление равно кавитационному~($p = p_{\mathrm{cav}}$) и~поток массы через границу непрерывен (баланс~\eqref{eq:jfo_couette_balance} из~раздела~3.5). Именно второе условие является дополнительным по~сравнению с~упрощёнными моделями и~обеспечивает выполнение баланса масс. Численно свободная граница отслеживается неявно через обновление активного множества~-- алгоритм, описанный в~разделах~2.3.3 и~3.5,~-- без явного хранения координат линии раздела зон~\cite{hori2002,szeri1998}.
